\documentclass[11pt,a4paper]{article}
\usepackage[utf8]{inputenc}
\usepackage{amsmath,amssymb,amsfonts}
\usepackage{booktabs}
\usepackage{geometry}
\usepackage{hyperref}
\usepackage{microtype}

\geometry{margin=25mm}

\title{Supplementary Materials: Self-Organized cancer fragility emerges from chromosomal instability}

\author{\textbf{Theoretical Biology \& Mathematical Oncology Framework}}
\date{\today}

\begin{document}

\maketitle


\section{Stationary Mutational Population distributions in the Liquid case}

We consider a liquid tumor model with fixed population size $N$. The number of wild-type cells is indicated with $W$, the number of dead cells with $D$. The cancer cell population can be further divided in $N_{\mathcal{I}}$ sub-populations (one per instability class), with $N_{\mathcal{I}}$ the number of instability genes. We indicate the abundance of each cancer sub-population with $C_k$, where $k \in \{1, \dots, N_I\}$ counts the number of mutated instability genes. We work with densities (fractions) rather than absolute abundances:

\begin{equation}
    f_W \equiv \frac{W}{N}, \quad f_k \equiv \frac{C_k}{N}, \quad p_D \equiv \frac{D}{N}
\end{equation}

These are constrained by the fixed-size normalization:
\begin{equation}
    \sum_{k=1}^{N_I} f_k + f_W + p_D = 1
\end{equation}

\subsubsection*{Survival probability after a mutation round}

Both cancer and wild-type cells undergo replication followed by a stochastic mutation event. After mutation, a \emph{viability check} is performed on all $N_{HK}$ housekeeping genes: if \emph{any} housekeeping gene becomes fully mutated (on all chromosome copies), the cell dies. Each gene has an independent per-locus mutation probability $\mu_k = k \cdot \delta\mu$ for a class-$k$ cell. The probability that a \emph{single} housekeeping gene survives (i.e.\ is \textbf{not} mutated) in one round is $(1 - k\,\delta\mu)$. Since all $N_{HK}$ housekeeping genes must independently survive:

\begin{equation}
    P_{s,k} = \underbrace{(1 - k\,\delta\mu)}_{\text{one HK gene survives}} \times \cdots \times \underbrace{(1 - k\,\delta\mu)}_{N_{HK}\text{ times}} = (1 - k\,\delta\mu)^{N_{HK}}
\end{equation}

\subsubsection*{The post-replication fates (exact multi-step model)}

After a mutation round, a cell in class $k$ can either remain in class $k$ (cloning), promote to class $k+j$ with $j \ge 1$ (multi-step promotion), or die due to a housekeeping mutation. Since mutations at different loci occur independently with per-locus probability $\mu_k = k \cdot \delta\mu$, we derive these probabilities exactly:

\paragraph{1. Faithful cloning ($\alpha_k$):} The cell remains in class $k$. This requires that none of the $N_{HK}$ housekeeping genes and none of the $(N_I - k)$ remaining unmutated instability genes are hit. The total number of loci that must remain unmutated is $N_{HK} + N_I - k$:
\begin{equation}
    \alpha_k = (1 - k\,\delta\mu)^{N_{HK} + N_I - k}
\end{equation}

\paragraph{2. Promotion ($P_{k \to k+j}$):} The cell advances from class $k$ to class $k+j$ (where $j \in \{1, \dots, N_I - k\}$ instability genes mutate, while all other instability loci and all housekeeping loci remain unmutated). Since the selection of which instability genes mutate is combinatorial:
\begin{equation}
    P_{k \to k+j} = \binom{N_I - k}{j} (k\,\delta\mu)^j (1 - k\,\delta\mu)^{N_{HK} + N_I - k - j}
\end{equation}
Note that faithful cloning is simply the diagonal case $\alpha_k = P_{k \to k}$.

\paragraph{3. Lethal mutation ($\delta_k$):} The cell dies because at least one housekeeping gene is hit, which is the complement of the survival probability $P_{s,k} = (1 - k\,\delta\mu)^{N_{HK}}$:
\begin{equation}
    \delta_k = 1 - P_{s,k} = 1 - (1 - k\,\delta\mu)^{N_{HK}}
\end{equation}

\textbf{Exact normalization check:} The sum of all possible survival outcomes plus the death probability is exactly 1:
\begin{equation}
    \sum_{j=0}^{N_I - k} P_{k \to k+j} + \delta_k = (1 - k\,\delta\mu)^{N_{HK}} \sum_{j=0}^{N_I - k} \binom{N_I - k}{j} (k\,\delta\mu)^j (1 - k\,\delta\mu)^{N_I - k - j} + [1 - (1 - k\,\delta\mu)^{N_{HK}}] = 1
\end{equation}
This binomial formulation provides an exact description of the mutation process without needing leading-order approximations.

\subsubsection*{Moran competition probabilities}

When a dividing cell targets a living (non-dead) cell, the outcome is decided by a \emph{Moran process}: the divider replaces the target with probability proportional to its fitness relative to the combined fitness. All cancer classes share the same intrinsic division rate $r_{\text{max}}$, while wild-type cells divide at rate $r_0$:

\begin{itemize}
    \item \textbf{Cancer vs Wild-Type:} The cancer cell (rate $r_{\text{max}}$) competes against a wild-type cell (rate $r_0$):
    \begin{equation}
        \beta_{cw} = \frac{r_{\text{max}}}{r_{\text{max}} + r_0}
    \end{equation}
    \item \textbf{Cancer vs Cancer:} Two cancer cells with the same rate $r_{\text{max}}$ compete symmetrically:
    \begin{equation}
        \beta_{cc} = \frac{r_{\text{max}}}{r_{\text{max}} + r_{\text{max}}} = \frac{1}{2}
    \end{equation}
    \item \textbf{Wild-Type vs Cancer:} The complement of the cancer-vs-WT outcome:
    \begin{equation}
        \beta_{wc} = \frac{r_0}{r_{\text{max}} + r_0} = 1 - \beta_{cw}
    \end{equation}
\end{itemize}

System parameters are formalized in Table~\ref{tab:parameters}.

\begin{table}[htbp]
\centering
\caption{Systematic Overview of the Unified Parameter Space}
\label{tab:parameters}
\begin{tabular}{lll}
\toprule
\textbf{Symbol} & \textbf{Biological Definition} & \textbf{Functional Form / Constraint} \\
\midrule
$N$ & Total fixed tissue lattice sites & Constant volume constraint \\
$N_I$ & Maximum viable cancer mutational classes & Total mutable Instability genes \\
$N_{HK}$ & Total critical Housekeeping genes & Continuous fitness check loci \\
$\delta\mu$ & Per-locus mutation rate increment & Constant scaling factor \\
$r_{\text{max}}$ & Intrinsic division rate of class $k$ cancer & Phenotypic fitness parameter \\
$r_0$ & Active division rate of healthy wild-type & Homeostatic driver baseline \\
$\alpha_k$ & Probability of a faithful clone & $(1-k\,\delta\mu)^{N_{HK} + N_I - k}$ \\
$P_{k \to k+j}$ & Probability of promotion ($k \to k+j$) & $\binom{N_I-k}{j}(k\,\delta\mu)^j(1-k\,\delta\mu)^{N_{HK}+N_I-k-j}$ \\
$\delta_k$ & Lethal mutation probability & $1 - (1-k\,\delta\mu)^{N_{HK}}$ \\
$\beta_{cw}$ & Moran win probability of cancer vs WT & $r_{\text{max}} / (r_{\text{max}} + r_0)$ \\
$\beta_{cc}$ & Moran win probability of cancer $k$ vs $j$ & $r_{\text{max}} / (r_{\text{max}} + r_{\text{max}})$ \\
$\beta_{wc}$ & Moran win probability of WT vs cancer $k$ & $r_0 / (r_{\text{max}} + r_0) = 1 - \beta_{cw}$ \\
\bottomrule
\end{tabular}
\end{table}

\subsection{Microscopic Mechanisms and Elementary Transition Rates}

After each duplication step, both mother and daughter cells independently undergo stochastic mutation events as described above. Both cells independently draw from the same probability distribution: cloning ($\alpha_k$), promoting ($P_{k \to k+j}$), or dying ($\delta_k$). Crucially, DNA replication and subsequent viability/mutation checks are only performed upon successful division (Moran competition success). If division is aborted, neither the mother nor the daughter replicates or mutates. Thus, all division-dependent rates (cloning, promotion, and death) for both mother and daughter cells are scaled by the division success probability $\phi_k$.

\subsubsection*{Effective reachable space $\phi_k$}

We define the fractional reachable space $\phi_k$ as the effective probability that a dividing class-$k$ cancer cell successfully places its daughter. The daughter can land on three types of targets:

\begin{equation}
    \phi_k = \underbrace{p_D}_{\substack{\text{dead slots:}\\\text{always available}}}
           + \underbrace{f_W\,\beta_{cw}}_{\substack{\text{WT cells displaced}\\\text{via Moran competition}}}
           + \underbrace{\sum_{j=1}^{N_I} f_j\,\beta_{cc}}_{\substack{\text{cancer cells displaced}\\\text{(50\% win rate)}}}
\end{equation}

\begin{itemize}
    \item $p_D$: the fraction of dead (empty) slots. Replacement always succeeds (no competition needed).
    \item $f_W \beta_{cw}$: wild-type cells are encountered with frequency $f_W$, and the cancer cell wins the Moran competition with probability $\beta_{cw}$.
    \item $\sum_j f_j \beta_{cc}$: other cancer cells are encountered with total frequency $\sum_j f_j$, with a symmetric 50\% win rate $\beta_{cc} = 1/2$.
\end{itemize}

\subsubsection*{Enumeration of all elementary transitions}

We now enumerate every possible microscopic event and its rate. The key insight is that during a single division event of a class-$k$ cancer cell, \emph{two independent fates} are resolved: (i) the \textbf{mother fate} (what happens to the original cell on its site after mutation), and (ii) the \textbf{daughter fate} (what happens to the copy placed at the target site after mutation). Both fates are conditioned on the division being successfully completed.

The complete transition rate matrix is given in Table~\ref{tab:transitions}. We explain each row below.

\begin{table}[htbp]
\centering
\caption{Complete Microscopic Transition Rate Matrix (Symmetrical Hazards)}
\label{tab:transitions}
\begin{tabular}{lllll}
\toprule
\textbf{Active Vector} & \textbf{Target \& Event Type} & $\mathbf{\Delta C_k}$ & $\mathbf{\Delta D}$ & \textbf{Transition Rate ($\lambda$)} \\
\midrule
Cancer Class $k$ & \textbf{Mother Fate:} Dies on Site & $-1$ & $+1$ & $C_k \, r_{\text{max}} \, \phi_k \, \delta_k$ \\
Cancer Class $k$ & \textbf{Mother Fate:} Promotes on Site ($k \to k+j$) & $-1_{k}, +1_{k+j}$ & $0$ & $C_k \, r_{\text{max}} \, \phi_k \, P_{k \to k+j}$ \\
Cancer Class $k$ & \textbf{Daughter Fate:} Clones on $D$ & $+1$ & $-1$ & $C_k \, r_{\text{max}} \, p_D \, \alpha_k$ \\
Cancer Class $k$ & \textbf{Daughter Fate:} Promotes on $D$ ($k \to k+j$) & $+1_{k+j}$ & $-1$ & $C_k \, r_{\text{max}} \, p_D \, P_{k \to k+j}$ \\
Cancer Class $k$ & \textbf{Daughter Fate:} Dies on $D$ & $0$ & $0$ & $C_k \, r_{\text{max}} \, p_D \, \delta_k$ \\
Cancer Class $k$ & \textbf{Daughter Fate:} Clones on $W$ & $+1$ & $0$ & $C_k \, r_{\text{max}} \, f_W \, \beta_{cw} \, \alpha_k$ \\
Cancer Class $k$ & \textbf{Daughter Fate:} Dies on $W$ & $0$ & $+1$ & $C_k \, r_{\text{max}} \, f_W \, \beta_{cw} \, \delta_k$ \\
Wild-Type ($W$) & WT Divides, Replaces Dead Void & $0$ & $-1$ & $W \, r_0 \, p_D$ \\
Wild-Type ($W$) & WT Divides, Clears Cancer $k$ & $-1$ & $0$ & $W \, r_0 \, f_k \, \beta_{wc}$ \\
\bottomrule
\end{tabular}
\end{table}

\paragraph{Explanation of each transition rate:}

\begin{enumerate}
    \item \textbf{Mother dies on site} ($C_k\, r_{\text{max}}\, \phi_k\, \delta_k$): A class-$k$ cell successfully initiates division (probability $\phi_k$), and the mother undergoes a lethal housekeeping mutation (probability $\delta_k$). The cell dies in place: $C_k \to C_k - 1$, $D \to D + 1$.

    \item \textbf{Mother promotes on site} ($C_k\, r_{\text{max}}\, \phi_k\, P_{k \to k+j}$): The mother successfully initiates division and mutates $j$ instability genes, advancing to class $k+j$ \emph{without leaving its site}. This is a type-swap: $C_k \to C_k - 1$, $C_{k+j} \to C_{k+j} + 1$, with no change to $D$.

    \item \textbf{Daughter clones on dead slot} ($C_k\, r_{\text{max}}\, p_D\, \alpha_k$): The daughter targets a dead slot (probability $p_D$) and faithfully clones (probability $\alpha_k$). Net effect: $C_k \to C_k + 1$, $D \to D - 1$.

    \item \textbf{Daughter promotes on dead slot} ($C_k\, r_{\text{max}}\, p_D\, P_{k \to k+j}$): The daughter targets a dead slot and mutates $j$ instability genes during division. Net effect: $C_{k+j} \to C_{k+j} + 1$, $D \to D - 1$.

    \item \textbf{Daughter dies on dead slot} ($C_k\, r_{\text{max}}\, p_D\, \delta_k$): The daughter targets a dead slot but suffers a lethal mutation. The dead slot remains dead, and no net population change occurs ($\Delta C_k = 0$, $\Delta D = 0$).

    \item \textbf{Daughter clones on WT} ($C_k\, r_{\text{max}}\, f_W\, \beta_{cw}\, \alpha_k$): The daughter targets a wild-type cell (probability $f_W$), wins the Moran competition (probability $\beta_{cw}$), and faithfully clones ($\alpha_k$). The WT cell is displaced: $C_k \to C_k + 1$, $W \to W - 1$.

    \item \textbf{Daughter dies on WT} ($C_k\, r_{\text{max}}\, f_W\, \beta_{cw}\, \delta_k$): The daughter wins competition against a WT cell but then suffers a lethal mutation. The WT cell is lost and a dead cell appears: $D \to D + 1$, $W \to W - 1$.

    \item \textbf{WT replaces dead slot} ($W\, r_0\, p_D$): A wild-type cell divides (rate $r_0$) and fills a dead slot (probability $p_D$). Net: $W \to W + 1$, $D \to D - 1$. No Moran competition is needed against dead cells.

    \item \textbf{WT clears cancer} ($W\, r_0\, f_k\, \beta_{wc}$): A wild-type cell targets a class-$k$ cancer cell (probability $f_k$) and wins the Moran competition with probability $\beta_{wc}$. Net: $C_k \to C_k - 1$, $W \to W + 1$.
\end{enumerate}

\subsection{Mean-Field Ordinary Differential Equations}

By taking the thermodynamic limit ($N \to \infty$), the stochastic master equation converges to a deterministic system of ODEs. We derive the ODE for each compartment by summing all transition rates from Table~\ref{tab:transitions} that affect that compartment, converting from absolute counts to densities ($f_k = C_k/N$).

\subsubsection*{Derivation of $df_k/dt$ for interior cancer classes ($1 < k < N_I$)}

For an interior cancer class $k$ ($1 < k < N_I$), the rate of change of the cell fraction $f_k$ is determined by five primary biological processes: (i) net clonal growth and survival ($r_{\text{max}} f_k \phi_k (2\alpha_k - 1)$), which balances faithful daughter cloning against mother cell loss via death or promotion; (ii) self-displacement loss within the same class ($r_{\text{max}} f_k^2 / 2$); (iii) cross-class loss due to displacement by other cancer classes ($f_k \sum_{j \neq k} r_{\text{max}} f_j \beta_{cc}$); (iv) wild-type clearance loss ($r_0 f_W f_k \beta_{wc}$); and (v) promotion influx from all upstream classes $i < k$ ($\sum_{i=1}^{k-1} 2\, r_{\text{max}} f_i \phi_i P_{i \to k}$). Combining these contributions yields:
\begin{equation}
    \boxed{\frac{df_k}{dt} = r_{\text{max}} f_k \phi_k (2\alpha_k - 1)
    - r_{\text{max}} \frac{f_k^2}{2}
    - f_k \!\sum_{j \neq k}\! r_{\text{max}} f_j \beta_{cc}
    - r_0 f_W f_k \beta_{wc}
    + \sum_{i=1}^{k-1} 2\, r_{\text{max}} f_i \phi_i P_{i \to k}}
\end{equation}

\subsubsection*{Boundary class $k = 1$}

The founding cancer class $k=1$ receives \textbf{no upstream mutational influx} (there is no class $k=0$). The equation is identical to the interior case but without the promotion influx:
\begin{equation}
    \frac{df_1}{dt} = r_{\text{max}} f_1 \phi_1 (2\alpha_1 - 1) - r_{\text{max}} \frac{f_1^2}{2} - f_1 \sum_{j \neq 1} r_{\text{max}} f_j \beta_{cc} - r_0 f_W f_1 \beta_{wc}
\end{equation}

\subsubsection*{Terminal class $k = N_I$}

For the terminal mutational class, all instability genes are already mutated, so further promotion is impossible: $P_{N_I \to N_I+j} = 0$. The mother-loss term simplifies to $-r_{\text{max}} f_{N_I} \phi_{N_I} (1 - \alpha_{N_I})$ where $1 - \alpha_{N_I} = \delta_{N_I}$. The promotion influx from all classes $i < N_I$ is present:
\begin{equation}
    \frac{df_{N_I}}{dt} = r_{\text{max}} f_{N_I} \phi_{N_I} (2\alpha_{N_I} - 1) - r_{\text{max}} \frac{f_{N_I}^2}{2} - f_{N_I} \sum_{j \neq N_I} r_{\text{max}} f_j \beta_{cc} - r_0 f_W f_{N_I} \beta_{wc} + \sum_{i=1}^{N_I-1} 2\, r_{\text{max}} f_i \phi_i P_{i \to N_I}
\end{equation}

\subsubsection*{Wild-type compartment $f_W$}

The WT population grows when WT cells divide and successfully place daughters on dead slots or displace cancer cells, and shrinks when cancer daughters displace WT cells:

\begin{itemize}
    \item \textbf{Gain:} WT divides onto dead slots (rate $r_0 f_W p_D$) or displaces cancer class $j$ (rate $r_0 f_W f_j \beta_{wc}$ per class).
    \item \textbf{Loss:} Cancer class $k$ daughters displace WT cells (rate $r_{\text{max}} f_k f_W \beta_{cw}$ per class).
\end{itemize}

\begin{equation}
    \frac{df_W}{dt} = \underbrace{r_0 f_W \left( p_D + \sum_{j=1}^{N_I} f_j \beta_{wc} \right)}_{\text{WT expansion (dead slots + cancer clearance)}} - \underbrace{f_W \sum_{k=1}^{N_I} r_{\text{max}} f_k \beta_{cw}}_{\text{displacement by cancer daughters}}
\end{equation}

\subsubsection*{Dead compartment $p_D$}

Dead cells accumulate from lethal mutations (both mother dying in place and daughter dying after placement) and are consumed when any living cell fills a dead slot. Since DNA replication only happens during successful division, both mother and daughter can only die if division succeeds.

\begin{itemize}
    \item \textbf{Gain from mother death:} $r_{\text{max}} f_k \phi_k \delta_k$ per class (row 1).
    \item \textbf{Gain from daughter death:} $r_{\text{max}} f_k \phi_k \delta_k$ per class (combining rows 5, 7, and cancer-cancer targets).
    \item \textbf{Total death-related gain:} $2\, r_{\text{max}} f_k \phi_k \delta_k$ per class.
    \item \textbf{Loss from filling:} Dead slots are consumed by cancer daughters ($r_{\text{max}} f_k p_D$, regardless of fate) and WT daughters ($r_0 f_W p_D$).
\end{itemize}

\begin{equation}
    \frac{dp_D}{dt} = \underbrace{\sum_{k=1}^{N_I} r_{\text{max}} f_k \left[ 2 \phi_k \delta_k - p_D \right]}_{\text{cancer-driven death accumulation minus filling}} - \underbrace{r_0 f_W p_D}_{\text{WT filling dead slots}}
\end{equation}

\subsection{Stationary State Distributions and Stability Analysis}

\subsubsection*{Linear stability of the healthy fixed point}

We examine the linear stability of the healthy fixed point $(f_W^* = 1,\; p_D^* = 0,\; f_k^* = 0\;\forall k)$ against an infinitesimal perturbation $f_1 = \epsilon \ll 1$, with all other cancer classes still at zero.

Starting from the $k=1$ boundary ODE and substituting the healthy fixed-point values, each term simplifies as follows:

\begin{itemize}
    \item \textbf{Mother loss \& daughter clone (combined):} $r_{\text{max}}\,\epsilon\,\phi_1\,(2\alpha_1 - 1)$. At the fixed point, $\phi_1 \to f_W \beta_{cw} = 1 \cdot \beta_{cw}$. This gives $r_{\text{max}}\,\beta_{cw}\,(2\alpha_1 - 1)\,\epsilon$.
    \item \textbf{Self-displacement:} $-r_{\text{max}}\,\epsilon^2/2 \approx 0$ --- vanishes at $\mathcal{O}(\epsilon^2)$.
    \item \textbf{Cross-class loss:} $-\epsilon \sum_{j\neq 1} r_{\text{max}} f_j \beta_{cc} = 0$.
    \item \textbf{WT clearance:} $-r_0\,(1)\,\epsilon\,\beta_{wc} = -r_0\,\beta_{wc}\,\epsilon$.
    \item \textbf{Promotion influx:} absent for $k=1$.
\end{itemize}

Collecting all $\mathcal{O}(\epsilon)$ terms:
\begin{equation}
    \frac{df_1}{dt} \approx \left[ r_{\text{max}}\,\beta_{cw}\,(2\alpha_1 - 1) - r_0\,\beta_{wc} \right] \epsilon
\end{equation}

Substituting $\beta_{cw} = r_{\text{max}}/(r_{\text{max}} + r_0)$ and $\beta_{wc} = r_0/(r_{\text{max}} + r_0)$:
\begin{equation}
    \frac{df_1}{dt} \approx \frac{r_{\text{max}}}{r_{\text{max}} + r_0} \left[ r_{\text{max}} (2\alpha_1 - 1) - r_0 \right] f_1
\end{equation}

\subsubsection*{Derivation of the growth foothold inequality}

The tumor can invade when the bracketed growth rate is positive:
\begin{equation}
    r_{\text{max}} (2\alpha_1 - 1) - r_0 > 0
\end{equation}

This yields the \textbf{Unified Growth Foothold Inequality}:
\begin{equation}
    \boxed{r_{\text{max}} (2\alpha_1 - 1) > r_0}
\end{equation}

\textbf{Physical interpretation:} The left-hand side balances the cancer cell's maximum birth rate times its net clonal expansion factor $(2\alpha_1 - 1)$, which accounts for the survival and fidelity of both the mother and daughter during a division event. The right-hand side $r_0$ represents the WT homeostatic division pressure. When mutation rates are too high, the cancer's net growth rate drops below the wild-type's division rate, preventing the tumor from taking foothold.

\subsubsection*{Expression in terms of the survival probability $P_{s,1}$}

We can rewrite the inequality in terms of the survival probability $P_{s,1} = (1-\delta\mu)^{N_{HK}}$ and the probability of faithful cloning $\alpha_1 = P_{s,1} (1-\delta\mu)^{N_I-1}$:
\begin{equation}
    \boxed{r_{\text{max}} \Big[ 2 P_{s,1}\,(1-\delta\mu)^{N_I-1} - 1\Big] > r_0}
\end{equation}

This form makes the phase boundary explicit: the critical $\delta\mu^*$ is the value at which $P_{s,1}$ (which decreases with $\delta\mu$) drops low enough to violate the inequality.

\subsubsection*{Stationary state: reduction to sequential quadratic equations}

When the growth inequality is satisfied, the system reaches a tumor-dominated stationary state where $df_k/dt = 0$ for all $k$. Setting the interior ODE to zero and rearranging:

\begin{equation}
    0 = r_{\text{max}} f_k^* \phi_k^* (2\alpha_k - 1) - r_{\text{max}} \frac{(f_k^*)^2}{2} - f_k^* \sum_{j \neq k} r_{\text{max}} f_j^* \beta_{cc} - r_0 f_W^* f_k^* \beta_{wc} + \sum_{i=1}^{k-1} 2\, r_{\text{max}} f_i^* \phi_i^* P_{i \to k}
\end{equation}

We separate the $\phi_k^*$ sum into its self-interaction part ($f_k^* \beta_{cc}$) and the rest ($\phi_k^{(\neq k)*}$):
\begin{equation}
    \phi_k^* = \phi_k^{(\neq k)*} + f_k^*\,\beta_{cc}
\end{equation}

Substituting and grouping terms by powers of $f_k^*$:

\begin{align}
    0 &= r_{\text{max}} f_k^* \left[\phi_k^{(\neq k)*} + f_k^* \beta_{cc}\right] (2\alpha_k - 1) - r_{\text{max}} \frac{(f_k^*)^2}{2} \nonumber \\
      &\quad - f_k^* \sum_{j \neq k} r_{\text{max}} f_j^* \beta_{cc} - r_0 f_W^* f_k^* \beta_{wc} + \sum_{i=1}^{k-1} 2\, r_{\text{max}} f_i^* \phi_i^* P_{i \to k} \nonumber \\
    0 &= r_{\text{max}} f_k^* \phi_k^{(\neq k)*} (2\alpha_k - 1) + r_{\text{max}} (f_k^*)^2 \beta_{cc} (2\alpha_k - 1) - r_{\text{max}} \frac{(f_k^*)^2}{2} \nonumber \\
      &\quad - f_k^* \sum_{j \neq k} r_{\text{max}} f_j^* \beta_{cc} - r_0 f_W^* f_k^* \beta_{wc} + \sum_{i=1}^{k-1} 2\, r_{\text{max}} f_i^* \phi_i^* P_{i \to k}
\end{align}

Since $\beta_{cc} = 1/2$, the quadratic terms combine as:
\begin{equation}
    -r_{\text{max}} (1 - \alpha_k) (f_k^*)^2
\end{equation}

Collecting terms, this takes the standard quadratic form:

\begin{equation}
    A_k\,(f_k^*)^2 - B_k\,f_k^* - Q_{k-1}^* = 0
\end{equation}

where the coefficients are:

\paragraph{$A_k$ --- Quadratic (self-interaction) coefficient:}
\begin{equation}
    A_k = r_{\text{max}} (1 - \alpha_k)
\end{equation}
This arises from the combination of self-displacement ($\beta_{cc} = 1/2$) and the mutation outcomes.

\paragraph{$B_k$ --- Linear (net growth) coefficient:}
\begin{equation}
    B_k = r_{\text{max}}\,\phi_k^{(\neq k)*}\,(2\alpha_k - 1) - \sum_{j \neq k} r_{\text{max}} f_j^*\,\beta_{cc} - r_0\,f_W^*\,\beta_{wc}
\end{equation}
This coefficient represents the \emph{net per-capita growth rate} of class $k$ at the stationary state, excluding self-interaction and upstream promotion.

\paragraph{$Q_{k-1}^*$ --- Upstream promotion source:}
\begin{equation}
    Q_{k-1}^* = \sum_{i=1}^{k-1} 2\, r_{\text{max}}\,f_i^*\,\phi_i^*\,P_{i \to k}
\end{equation}
This is the rate at which cells flow into class $k$ from all lower classes $i < k$ (dual-sourced: mother promoting + daughter promoting). For $k=1$, $Q_0^* = 0$.

\subsubsection*{Solution via the quadratic formula}

The quadratic $A_k x^2 - B_k x - Q_{k-1}^* = 0$ (with $x = f_k^*$) has the physical (positive) root:
\begin{equation}
    \boxed{f_k^* = \frac{B_k + \sqrt{B_k^2 + 4\,A_k\,Q_{k-1}^*}}{2\,A_k}}
\end{equation}

The positive sign before the square root is chosen because $f_k^* \geq 0$ is required. The system is solved \emph{sequentially}: $f_1^*$ is determined first (with $Q_0^* = 0$), then $f_2^*$ using $f_1^*$, and so on up to $f_{N_I}^*$.

\end{document}